
\subsubsection{Dạng 1. Phân biệt các loại mô hình trong nghiên cứu Vật lí}
\begin{phuongphap}
\begin{enumerate}
\item Xác định đối tượng thực tế cần nghiên cứu và đặc điểm của mô hình được đưa ra trong đề bài.
\item Phân tích tính chất của mô hình: nếu là các thực thể vật chất trực quan thì thuộc mô hình vật chất; nếu dùng các khái niệm, giả thuyết trừu tượng thì thuộc mô hình lí thuyết; nếu sử dụng các biểu thức toán học, công thức hay đồ thị thì thuộc mô hình toán học.
\item Đối chiếu định nghĩa để phân loại chính xác mô hình tương ứng.
\end{enumerate}
\end{phuongphap}
\begin{vidumau}
Trong các mô hình sau đây được sử dụng trong học tập và nghiên cứu Vật lí: quả địa cầu thu nhỏ của Trái Đất, khái niệm chất điểm khi nghiên cứu chuyển động của ô tô chạy từ Hà Nội vào Thanh Hóa, công thức tính quãng đường trong chuyển động thẳng đều $s = v \cdot t$. Hãy phân loại các mô hình này vào các nhóm mô hình tương ứng.
\loigiai{
\begin{itemize}
\item Bước 1: Phân tích đối tượng và đặc điểm của từng mô hình được đề cập trong bài toán. Quả địa cầu là vật thu nhỏ mô phỏng Trái Đất thực tế. Khái niệm chất điểm thay thế cho ô tô có kích thước rất nhỏ so với quãng đường. Công thức $s = v \cdot t$ biểu diễn mối quan hệ định lượng giữa quãng đường, vận tốc và thời gian.
\item Bước 2: Áp dụng định nghĩa các loại mô hình để phân loại. Quả địa cầu là một vật thể thực, trực quan thay thế cho vật thật nên thuộc nhóm mô hình vật chất. Khái niệm chất điểm là một giả định lí thuyết bỏ qua kích thước của vật để đơn giản hóa việc tính toán chuyển động nên thuộc nhóm mô hình lí thuyết. Công thức $s = v \cdot t$ sử dụng toán học để mô tả quy luật vật lí nên thuộc nhóm mô hình toán học.
\item Bước 3: Kết luận phân loại chính xác các mô hình tương ứng: Quả địa cầu thuộc mô hình vật chất; khái niệm chất điểm thuộc mô hình lí thuyết; công thức $s = v \cdot t$ thuộc mô hình toán học.
\end{itemize}
}
\end{vidumau}

\subsubsection{Dạng 2. Vai trò của Vật lí đối với khoa học, kĩ thuật, công nghệ và đời sống}
\begin{phuongphap}
\begin{enumerate}
\item Đọc kĩ nội dung ứng dụng thực tế hoặc công nghệ được đề cập trong câu hỏi.
\item Liên hệ kiến thức lịch sử phát triển vật lí và bốn cuộc cách mạng công nghiệp để xác định lĩnh vực nghiên cứu vật lí cốt lõi đứng sau ứng dụng đó.
\item Đánh giá tác động đa chiều (tích cực như nâng cao năng suất, tiết kiệm năng lượng; hoặc hạn chế như gây ô nhiễm, biến đổi khí hậu) của công nghệ đó đối với đời sống và môi trường.
\end{enumerate}
\end{phuongphap}
\begin{vidumau}
Hãy phân tích vai trò của việc phát minh ra động cơ hơi nước đối với sự phát triển của xã hội loài người và cho biết phát minh này dựa trên phân ngành nào của Vật lí học.
\loigiai{
\begin{itemize}
\item Bước 1: Phân tích đối tượng nghiên cứu. Động cơ hơi nước là thiết bị chuyển hóa nhiệt năng của hơi nước thành cơ năng để vận hành máy móc sản xuất và phương tiện giao thông như tàu hỏa, tàu thủy.
\item Bước 2: Xác định cơ sở vật lí và vai trò lịch sử. Động cơ hơi nước hoạt động dựa trên các quy luật của Nhiệt học và Nhiệt động lực học. Sự ra đời của động cơ hơi nước đã thay thế sức lao động thủ công của con người và sức kéo động vật bằng sức máy, mở đầu cho cuộc cách mạng công nghiệp lần thứ nhất ở thế kỉ XVIII, làm thay đổi toàn diện phương thức sản xuất của nhân loại.
\item Bước 3: Kết luận vai trò và cơ sở khoa học. Phát minh động cơ hơi nước dựa trên nền tảng của Nhiệt động lực học, đóng vai trò quyết định thúc đẩy sự ra đời của nền sản xuất cơ khí hóa quy mô lớn và cuộc cách mạng công nghiệp lần thứ nhất.
\end{itemize}
}
\end{vidumau}

\subsubsection{Dạng 3. Vận dụng tiến trình của phương pháp thực nghiệm}
\begin{phuongphap}
\begin{enumerate}
\item Xác định rõ mục đích nghiên cứu, hiện tượng hoặc quy luật vật lí cần được kiểm chứng thông qua thực nghiệm.
\item Sắp xếp và phân tích các bước trong tiến trình thực nghiệm: từ việc đề xuất giả thuyết ban đầu, thiết kế và bố trí dụng cụ thí nghiệm, tiến hành đo đạc thu thập dữ liệu, đến xử lí sai số và kiểm chứng giả thuyết.
\item Nhận định tính đúng đắn của giả thuyết hoặc đưa ra quy luật tổng quát từ dữ liệu thực nghiệm đã thu thập được.
\end{enumerate}
\end{phuongphap}
\begin{vidumau}
Trình bày phương pháp thực nghiệm của Galileo Galilei khi tiến hành nghiên cứu sự rơi của các vật để bác bỏ quan niệm của Aristotle rằng vật nặng rơi nhanh hơn vật nhẹ.
\loigiai{
\begin{itemize}
\item Bước 1: Xác định mâu thuẫn lí thuyết và mục đích thực nghiệm. Aristotle cho rằng lực cản và trọng lượng quyết định tốc độ rơi của vật, do đó vật nặng phải rơi nhanh hơn vật nhẹ. Galileo muốn kiểm chứng quan điểm này bằng cách thực hiện thí nghiệm trực tiếp trong thực tế.
\item Bước 2: Phân tích quá trình thực nghiệm của Galileo. Ông đã chọn tháp nghiêng Pisa làm nơi thực hiện thí nghiệm để đảm bảo độ cao đủ lớn cho việc quan sát. Ông cho thả rơi đồng thời hai quả cầu bằng chất liệu giống nhau nhưng có khối lượng chênh lệch nhau rất nhiều từ trên đỉnh tháp xuống đất. Kết quả thực tế cho thấy hai quả cầu chạm đất gần như cùng một lúc.
\item Bước 3: Kết luận ý nghĩa thực nghiệm. Kết quả thực nghiệm của Galileo đã trực tiếp bác bỏ quan niệm suy luận chủ quan của Aristotle, khẳng định rằng khi loại bỏ hoặc giảm thiểu tối đa sức cản của không khí thì các vật nặng nhẹ khác nhau đều rơi nhanh như nhau, đặt nền móng cho sự ra đời của phương pháp thực nghiệm trong Vật lí học.
\end{itemize}
}
\end{vidumau}

\subsubsection{Dạng 4. Vận dụng tiến trình của phương pháp mô hình}
\begin{phuongphap}
\begin{enumerate}
\item Xác định đối tượng thực tế phức tạp cần nghiên cứu trong bài toán cụ thể.
\item Lựa chọn các đặc tính quan trọng có ảnh hưởng lớn đến hiện tượng cần khảo sát, đồng thời lược bỏ các chi tiết phụ không ảnh hưởng nhiều nhằm đơn giản hóa đối tượng.
\item Xây dựng mô hình tương ứng (vật chất, lí thuyết hoặc toán học) để khảo sát, từ đó áp dụng các quy luật vật lí lên mô hình đó nhằm giải thích cho đối tượng thực tế ban đầu.
\end{enumerate}
\end{phuongphap}
\begin{vidumau}
Khi nghiên cứu chuyển động của một chiếc tàu thủy di chuyển trên chặng đường dài từ cảng Hải Phòng đến cảng Sài Gòn, người ta đã coi tàu thủy là một chất điểm. Hãy giải thích cơ sở vật lí của việc sử dụng mô hình lí thuyết này.
\loigiai{
\begin{itemize}
\item Bước 1: Xác định đối tượng thực tế và mô hình thay thế. Đối tượng thực tế là chiếc tàu thủy có kích thước chiều dài khoảng vài chục mét đến trăm mét chuyển động trên biển. Mô hình thay thế là chất điểm (một điểm có khối lượng bằng khối lượng của vật).
\item Bước 2: So sánh kích thước của vật với phạm vi chuyển động. Quãng đường từ Hải Phòng đến Sài Gòn dài khoảng hơn $1500\,\mathrm{km}$ (tương đương $1500000\,\mathrm{m}$), lớn hơn rất nhiều lần so với kích thước hình học của tàu thủy (khoảng vài chục mét). Do đó, hình dáng và kích thước của tàu thủy lúc này không ảnh hưởng đến việc xác định vị trí và quỹ đạo chuyển động tổng thể của nó trên bản đồ đường biển.
\item Bước 3: Kết luận cơ sở của phương pháp mô hình. Việc coi tàu thủy là một chất điểm là hoàn toàn hợp lí vì kích thước của tàu rất nhỏ so với quãng đường di chuyển, giúp đơn giản hóa bài toán động học mà vẫn đảm bảo độ chính xác của kết quả nghiên cứu vị trí và thời gian hành trình.
\end{itemize}
}
\end{vidumau}

\subsubsection{Dạng 5. Vận dụng tiến trình nghiên cứu Vật lí vào tình huống thực tiễn}
\begin{phuongphap}
\begin{enumerate}
\item Quan sát thế giới tự nhiên xung quanh để phát hiện các câu hỏi, hiện tượng vật lí cần giải thích.
\item Xây dựng quy trình nghiên cứu khoa học gồm 5 bước tiêu chuẩn: Đề xuất vấn đề, đưa ra giả thuyết, lập phương án kiểm chứng (lí thuyết hoặc thực nghiệm), tiến hành kiểm chứng và phân tích dữ liệu, rút ra kết luận khoa học.
\item Áp dụng quy trình này để thiết lập sơ đồ nghiên cứu cho một hiện tượng thực tế cụ thể.
\end{enumerate}
\end{phuongphap}
\begin{vidumau}
Một học sinh nhận thấy khi thả một tờ giấy phẳng và một tờ giấy được vo tròn có cùng khối lượng từ cùng một độ cao trong lớp học, tờ giấy vo tròn rơi xuống đất nhanh hơn rất nhiều so với tờ giấy phẳng. Hãy thiết lập tiến trình nghiên cứu gồm 5 bước để giải thích hiện tượng thực tế này.
\loigiai{
\begin{itemize}
\item Bước 1: Nhận diện hiện tượng và đề xuất vấn đề nghiên cứu. Quan sát thấy hai tờ giấy nặng như nhau nhưng hình dạng khác nhau lại rơi nhanh chậm khác nhau trong phòng học. Đặt câu hỏi: Phải chăng hình dạng của vật ảnh hưởng đến tốc độ rơi của vật trong không khí?
\item Bước 2: Đề xuất giả thuyết khoa học và phương án kiểm chứng. Giả thuyết: Tờ giấy phẳng rơi chậm hơn do diện tích tiếp xúc với không khí lớn hơn, chịu lực cản của không khí lớn hơn tờ giấy vo tròn. Thiết kế phương án thí nghiệm: Chuẩn bị hai tờ giấy giống hệt nhau, một tờ giữ phẳng, một tờ vo tròn chặt. Thả chúng rơi đồng thời ở cùng một độ cao và dùng camera quay chậm để đếm thời gian rơi.
\item Bước 3: Thực hiện kiểm chứng và kết luận. Tiến hành thả rơi và ghi nhận thời gian rơi của tờ giấy vo tròn ngắn hơn tờ giấy phẳng rất nhiều. Phân tích kết quả: Tờ giấy phẳng có bề mặt rộng hứng nhiều không khí hơn khi rơi, lực cản không khí tác dụng lên nó lớn hơn trọng lực nhiều so với tờ giấy vo tròn có diện tích bề mặt nhỏ. Rút ra kết luận: Sự rơi nhanh hay chậm của các vật trong không khí phụ thuộc vào độ lớn của lực cản không khí tác dụng lên vật, lực cản càng nhỏ so với trọng lực thì vật rơi càng nhanh.
\end{itemize}
}
\end{vidumau}

\subsubsection{Dạng 6. Đối tượng nghiên cứu và mục tiêu của môn Vật lí}
\begin{phuongphap}
\begin{enumerate}
\item Xác định hệ vật lý hoặc hiện tượng được nêu trong câu hỏi (ví dụ: dòng điện, ánh sáng, nguyên tử, vũ trụ).
\item Phân tích và phân loại đối tượng đó vào các phân ngành tương ứng của Vật lí học (Cơ học, Nhiệt học, Điện học, Từ học, Quang học, Vật lí hạt nhân).
\item Xác định xem đối tượng đó thuộc cấp độ nghiên cứu vi mô (thế giới các hạt cực nhỏ) hay vĩ mô (thế giới các vật thể thông thường và thiên văn học) để làm rõ mục tiêu tìm kiếm quy luật vận động tổng quát của nó.
\end{enumerate}
\end{phuongphap}
\begin{vidumau}
Xét hai đối tượng nghiên cứu sau của Vật lí học: tương tác giữa các proton và neutron bên trong hạt nhân nguyên tử heli, và sự chuyển động quỹ đạo của Trái Đất xung quanh Mặt Trời. Hãy xác định cấp độ nghiên cứu (vi mô hay vĩ mô) và phân ngành vật lí tương ứng cho từng đối tượng.
\loigiai{
\begin{itemize}
\item Bước 1: Phân tích đối tượng thứ nhất. Tương tác giữa proton và neutron diễn ra bên trong hạt nhân nguyên tử heli, có kích thước cực kì nhỏ bé (khoảng $10^{-15}\,\mathrm{m}$). Đối tượng này thuộc cấp độ nghiên cứu vi mô, liên quan đến cấu trúc sâu sắc nhất của vật chất và thuộc phân ngành Vật lí hạt nhân.
\item Bước 2: Phân tích đối tượng thứ hai. Sự chuyển động quỹ đạo của Trái Đất xung quanh Mặt Trời liên quan đến các thiên thể lớn trong không gian vũ trụ, có khoảng cách và kích thước hình học vô cùng lớn. Đối tượng này thuộc cấp độ nghiên cứu vĩ mô, liên quan đến lực hấp dẫn và thuộc phân ngành Cơ học thiên văn (Vũ trụ học).
\item Bước 3: Kết luận phân loại. Đối tượng tương tác trong hạt nhân heli thuộc cấp độ nghiên cứu vi mô (Vật lí hạt nhân). Đối tượng chuyển động của Trái Đất xung quanh Mặt Trời thuộc cấp độ nghiên cứu vĩ mô (Cơ học vũ trụ).
\end{itemize}
}
\end{vidumau}

\subsubsection{Dạng 7. Quá trình phát triển của Vật lí}
\begin{phuongphap}
\begin{enumerate}
\item Đọc kĩ câu hỏi để xác định mốc thời gian lịch sử, nhà bác học hoặc phát minh vật lí quan trọng được nhắc tới.
\item Đối chiếu phát kiến đó với ba giai đoạn phát triển chính của Vật lí học:
\begin{itemize}
\item Giai đoạn tiền Vật lí (từ thời cổ đại đến thế kỉ XVII): Quan sát tự nhiên và suy luận triết học, mang tính chủ quan.
\item Giai đoạn Vật lí cổ điển (thế kỉ XVII đến cuối thế kỉ XIX): Phát triển mạnh mẽ nhờ phương pháp thực nghiệm và các định luật Newton, Maxwell, Faraday.
\item Giai đoạn Vật lí hiện đại (thế kỉ XX đến nay): Ra đời thuyết tương đối của Einstein và cơ học lượng tử, mở rộng nghiên cứu sang cấp độ cực nhỏ và cực lớn.
\end{itemize}
\item Đánh giá ý nghĩa khoa học và triết học của phát minh đó đối với lịch sử khoa học tự nhiên.
\end{enumerate}
\end{phuongphap}
\begin{vidumau}
Sự ra đời của thuyết tương đối của Albert Einstein và thuyết lượng tử của Max Planck ở đầu thế kỉ XX đã đánh dấu bước chuyển sang giai đoạn phát triển nào của Vật lí học? Hãy trình bày ngắn gọn tầm quan trọng của hai thuyết này đối với hệ thống tri thức vật lí.
\loigiai{
\begin{itemize}
\item Bước 1: Xác định mốc thời gian và các học thuyết lớn. Cuối thế kỉ XIX và đầu thế kỉ XX, sự ra đời của thuyết lượng tử (Planck) và thuyết tương đối (Einstein) đã giải quyết những mâu thuẫn mà hệ thống Vật lí cổ điển của Newton không thể giải thích được ở cấp độ vi mô và vận tốc tiệm cận tốc độ ánh sáng.
\item Bước 2: Đối chiếu giai đoạn phát triển lịch sử. Bước ngoặt này đánh dấu sự chuyển mình từ Vật lí cổ điển sang giai đoạn Vật lí hiện đại. Thuyết tương đối của Einstein đã định nghĩa lại khái niệm không gian và thời gian ở vận tốc vũ trụ. Thuyết lượng tử của Planck đã mở cánh cửa nghiên cứu cấu trúc gián đoạn của năng lượng ở cấp độ nguyên tử và các hạt cơ bản.
\item Bước 3: Kết luận ý nghĩa lịch sử. Sự ra đời của hai thuyết này đánh dấu bước khởi đầu của giai đoạn Vật lí hiện đại, giúp loài người mở rộng biên giới tri thức từ thế giới vĩ mô thông thường xuống sâu thế giới vi mô hạt nhân và vươn xa tới quy mô toàn vũ trụ.
\end{itemize}
}
\end{vidumau}